\lecture{2}{Thu 23 Jan 2025 10:58}{hi}
I walked in and we're talking about some stuff with boats and physics? I don't really know what's happening, but we did some vector stuff to conclude the best way to sail against the wind is at a 30 degree angle.

It is usually too difficult to list all open sets in a given topology, so we use a basis instead.
\begin{definition}[Basis]\label{dfn:1.4}
	Let $X$ be a nonempty set and let $\mathcal{B} $ be a collection of subsets of $X$. Then $\mathcal{B} $ is a \emph{basis} of some topology $\mathcal{T}_{\mathcal{B} } $ on $X$ iff
	\begin{enumerate}
		\item For all $x\in X$, there exists $V\in \mathcal{B} $ such that $x\in V$;
		\item For all $V_1,V_2\in \mathcal{B} $ with $x\in V_1\cap V_2$, there exists $V_3\in \mathcal{B} $ such that $x\in V_3\subseteq V_1\cap V_2$.
	\end{enumerate}
\end{definition}
\begin{proposition}\label{prp:1.1}
	Let $\mathcal{B} $ be a basis for some set $X$. Then the set $\mathcal{T}_{\mathcal{B} }$ of the empty set, together with  all unions of elements of $\mathcal{B} $,  forms a topology on $X$.
\end{proposition}
\begin{proof}
	Recall the three axioms from definition \ref{dfn:1.1}. By definition, $\varnothing\in \mathcal{T}_\mathcal{B} $. By property (1) of definition \ref{dfn:1.4}, for all $x\in X$, there exists $V_x\in \mathcal{B} $ such that $x\in V_x$. By definition of $\mathcal{T} _\mathcal{B} $,
	\[
		X = \bigcup_{x\in X} V_x \in \mathcal{T}_{\mathcal{B} } 
	.\]
	Therefore, property (1) is satisfied.

	Now, we show property (2); closure under finite intersection. Let $\left\{ A_i \right\}_{i=1}^k \subseteq \mathcal{T} _{\mathcal{B} }$. By definition, we have that $A_i$ is either empty, or can be represented as a union
	\[
		A_i = \bigcup_{j\in J_i} V_{ij}
	\]
	for $\left\{ V_{ij} \right\}_{j\in J_i} \subseteq \mathcal{B} $. If any $A_i=\varnothing$, we have that the intersection
	\[
		\bigcap_{i=1}^{k} A_i=\varnothing\in \mathcal{T} _{\mathcal{B} }
	.\]
	Thus, we can assume WLOG that $A_i \neq \varnothing$ for all $i$. We have
	\[
		\bigcap_{i=1}^{k} A_i = \bigcap_{i=1}^{k} \bigcup_{j\in J_i} V_{ij}
	.\]
	Recall that the distributive property says $a \cup (b\cap c)=(a\cap b)\cup (a\cap c)$. We can thus rewrite something along the lines of
	\[
		\bigcap_{i=1}^{k} \bigcup_{j\in J_i} V_{ij}=\bigcup_{j_i\in J_i}\bigcap_{i\in I} V_{ij_i}
	.\]

	We thus need only prove this intersection is an element of $\mathcal{T}_{\mathcal{B} }$. We prove this by strong induction. For $k=1$, we have $V_1=\bigcap_{i=1}^{k} B_i\in \mathcal{B} \subseteq \mathcal{T} _{\mathcal{B} }$. Now assuming the inductive hypothesis for $n\leq k$, we have for any $\left\{ V_i \right\}_{i=1}^{k+1}\subseteq  \mathcal{B} $, the intersection
	\[
		\bigcap_{i=1}^{k} V_i \in \mathcal{B} 
	.\]
	We can thus represent it as a union of elements $\left\{ V_i' \right\}_{i\in I} \in \mathcal{T} _{\mathcal{B} }$. By distributivity, it follows that
	\[
		\bigcap_{i=1}^{k+1} V_{i}=\left( \bigcup_{i\in I} V_i'  \right) \cap V_{k+1}=\bigcup_{i\in I} \left( V_i' \cap V_{k+1} \right)  
	.\]
	By definition \ref{dfn:1.4}, for all $i$ and for all $x\in V_i'\cap V_{k+1}$, there exists $V_x \subseteq V_i'\cap V_{k+1}$ such that $x\in V_x$. We then have
	\[
		\bigcup_{i\in I}\left( V_i' \cap V_{k+1} \right) =\bigcup_{i\in I} \bigcup_{x\in V_i'\cap V_{k+1}} V_x 
	.\]
	Since $V_x\in \mathcal{B} $ for all $x$, by definition of $\mathcal{T} _{\mathcal{B} }$, we have
	\[
		\bigcup_{i\in I} \bigcup_{x\in V_i'\cap V_{k+1}} V_x \in \mathcal{T} _{\mathcal{B} }
	.\]
	By induction, we are done.

	Finally, we show property (3). Let $\left\{ A_i \right\}_{i\in I}\subseteq \mathcal{T} _\mathcal{B} $. By definition of $\mathcal{T}_{\mathcal{B} }$, for all $i$ there exists $\left\{ B_{ij} \right\}_{j\in J_i}$ such that
	\[
		A_i = \bigcup_{j\in J_i} B_{ij} 
	.\]
	We then have
	\[
		\bigcup_{i\in I} A_i = \bigcup_{i\in I}\bigcup_{j\in J_i} B_{ij}=\bigcup_{i\in I,\; j\in J_i} B_{ij}
	.\]
	Since this is a union of elements in $\mathcal{B} $, by definition of $\mathcal{T}_{\mathcal{B} } $, we have that
	\[
		\bigcup_{i\in I,\; j\in J_i} B_{ij}\in \mathcal{T} _{\mathcal{B} }
	.\]
\end{proof}
\begin{definition}[Topology Generated by a Basis]\label{dfn:1.5}
	Let $\mathcal{B} $ be a basis for a set $X$. Then the set $\mathcal{T} _\mathcal{B} $ consisting of the empty set together with all unions of elements in $\mathcal{B} $ is a topology on $X$, called the \emph{topology generated by $\mathcal{B} $}.
\end{definition}

For example, define 
\[
	B_r(x)\coloneqq \left\{ a\in\mathbb{R} \mid \left| x-a \right| <r \right\} 
.\]
Then the collection
\[
	\mathcal{B} =\left\{ B_r(x) \mid x\in \mathbb{R},r\in\mathbb{R}^+\right\} 
\]
is a basis for the standard topology on $\mathbb{R}$. This generalizes; let
\[
	B_r(x)\coloneqq \left\{ a\in\mathbb{R}^n \mid \left\lVert x-a \right\rVert <r \right\} 
.\]
Then
\[
	\mathcal{B} =\left\{ B_r(x) \mid r\in\mathbb{R}^+,x\in\mathbb{R}^n \right\} 
\]
is a basis for $\mathbb{R}^n$. We will prove this. Let $x\in \mathbb{R}^n$. Note that for any choice $r\in\mathbb{R}^+$, we have $x\in B_r(x)$. This shows property (1) of definition \ref{dfn:1.4}. As for property (2), let $x\in B_{r_2}(x_1)\cap B_{r_2}(x_2)$. It follows that this intersection is nonempty, which implies that $\left\lVert x_1-x \right\rVert ,\left\lVert x_2-x \right\rVert <r_1,r_2$.  We wish to find some $r$ such that $B_r(x)\subseteq B_{r_1}(x_1)\cap B_{r_2}(x_2)$. I propose we choose $\left| r_1 -r_2\right|$. Let $y\in B_{\left| r_1-r_2 \right| }(x)$. We wish to show $\left\lVert x_1-y \right\rVert <r_1$ and $\left\lVert x_2-y \right\rVert <r_2$. We have by the triangle inequality that
\[
	\left\lVert y-x_1 \right\rVert \leq \left\lVert y-x \right\rVert +\left\lVert x-x_1 \right\rVert 
.\]
We thus need only show this for $\left\lVert y-x_1 \right\rVert =\left\lVert y-x \right\rVert+\left\lVert x-x_1 \right\rVert  $. Note that
\[
	\left\lVert y-x \right\rVert <\left| r_1-r_2 \right| ,\qquad \left\lVert x-x_1 \right\rVert <r_1
.\]
We thus have
\[
	\left\lVert y-x_1 \right\rVert <\left| r_1-r_2 \right| +r_1=\left| r_2-r_1+r_1 \right| =r_1
,\]
because $r_1>0$. Therefore, $\left\lVert y-x_1 \right\rVert <r_1$. The proof is equivalent for $r_2$. Therefore, $B_{\left| r_1-r_2 \right| }(x)\subseteq B_{r_1}(x_1)\cap B_{r_2}(x_2)$, and thus by definition \ref{dfn:1.4}, $\mathcal{B} $ is a basis for $\mathbb{R}^n$.

Note that the only nontrivial thing used here is the triangle inequality. In the future, we will define a \emph{metric}, which is an abstraction of distance which maintains this, and a few other important properties of distance. We will then find that the collection of all open balls where the norm is replaced by the metric forms a basis for the metric space. This will be called the metric topology.

Note that the basis of a topology is not unique.
